\documentclass[conference]{IEEEtran}
\usepackage{graphicx}
\usepackage{url}
\usepackage{amsmath}
\usepackage{hyperref}
\usepackage{float}
\usepackage{booktabs}
\usepackage{multirow}
\usepackage{placeins}
\usepackage{verbatim}
\bibliographystyle{IEEEtran}

% ============================================================================
% MSDF-India -- SUPPLEMENTARY MATERIAL.
%
% Every count, shape, dtype, field name and file name in this document was read
% directly from the released Drive tree
%   https://drive.google.com/drive/folders/1bb0pH2RN5EvIYG701BpX3LoiHa3vy1QQ
% on 14 August 2026 (24 metadata.csv files, 2 summary_features.csv files and one
% clip's six .npy arrays were downloaded and inspected).
%
% Spots marked "% TODO(authors):" cannot be recovered from the release itself and
% must be filled from the collection / generation logs. They are NOT guesses.
% ============================================================================

\title{Supplementary Material\\
\large MSDF-India: A Multilingual Singing Deepfake Dataset\\for Cross-Lingual Analysis}

\begin{document}
\maketitle

% ============================================================================
\section{Scope}
% ============================================================================
This document accompanies the MSDF-India paper and specifies the released resource in
full: the complete metadata of the corpus, the exact feature-extraction settings, the
on-disk file structure and naming conventions, the mapping from feature files back to
source recordings, and the procedure for reproducing the cross-lingual benchmark. It is
intended to be sufficient for an independent group to reload the release, rebuild the
evaluation splits, and obtain comparable numbers without contacting the authors.

All counts, array shapes, data types and field names reported here were read from the
released archive itself rather than from the generation scripts, so that the document
describes what a user actually downloads. Where a property of the data cannot be
recovered from the release (for example, the per-recording source URL, or the target
speaker used for a particular conversion), this is stated explicitly rather than
inferred.

% ============================================================================
\section{What is Downloadable}
% ============================================================================
The release is distributed as a single Google Drive folder:
\begin{center}
\url{https://drive.google.com/drive/folders/1bb0pH2RN5EvIYG701BpX3LoiHa3vy1QQ}
\end{center}

\noindent It contains two top-level trees:

\begin{itemize}
  \item \texttt{Real\_dataset/} --- the bonafide partition, one folder per language.
  \item \texttt{Deepfake\_dataset/} --- the synthetic partition, one folder per
        generation method, each with one folder per language.
\end{itemize}

Every language folder, bonafide or synthetic, has the same internal layout: six
feature directories holding one \texttt{.npy} array per 10\,s clip, plus two CSV files
(\texttt{metadata.csv} and \texttt{summary\_features.csv}) describing every clip in that
folder. No audio waveforms are distributed; the release is feature-level only, and the
released arrays are non-invertible derived representations.

The archive holds \textbf{58{,}597 clips} of exactly 10.0\,s each --- 162.8\,h of
material --- comprising 18{,}086 bonafide clips (50.2\,h) and 40{,}511 synthetic clips
(112.5\,h). Per-clip storage is $\approx$1.07\,MiB across the six feature arrays, giving
a total download of $\approx$61\,GiB; the frame-level HNR arrays alone account for 57\%
of that volume (Section~\ref{sec:featspec}).

% ============================================================================
\section{Dataset Metadata}
% ============================================================================

\subsection{Bonafide Source Recordings}
Table~\ref{tab:s1} lists the bonafide corpus per language: the number of source
recordings admitted by the selection criteria, their total pre-segmentation duration, and
the number of fixed 10\,s clips they yield. Segmentation is deterministic
($\lfloor \text{duration}/10 \rfloor$ clips per recording, trailing fragment discarded),
so the clip counts follow from the source durations. One Punjabi recording of 4.9\,s is
removed by the minimum-duration rule, which is why 71 Punjabi source recordings appear as
70 in the released clips.

\begin{table}[t]
\centering
\caption{Bonafide partition: source recordings, pre-segmentation duration, released
10\,s clips and audio retention, per language.}
\label{tab:s1}
\footnotesize
\setlength{\tabcolsep}{4pt}
\begin{tabular}{lrrrrr}
\toprule
\textbf{Language} & \textbf{Songs} & \textbf{Src (min)} & \textbf{Clips} & \textbf{Clip h} & \textbf{Ret.} \\
\midrule
Hindi   & 84  & 1111.7 & 6{,}634 & 18.43 & 99.5\% \\
Punjabi & 71  &  404.7 & 2{,}393 &  6.65 & 98.6\% \\
Bengali & 64  &  473.0 & 2{,}806 &  7.79 & 98.9\% \\
Telugu  & 50  &  752.8 & 4{,}489 & 12.47 & 99.4\% \\
Odia    & 15  &  251.8 & 1{,}505 &  4.18 & 99.6\% \\
Tamil   & 10  &   44.0 &    259 &  0.72 & 98.0\% \\
\midrule
\textbf{Total} & \textbf{294} & \textbf{3038.0} & \textbf{18{,}086} & \textbf{50.24} & \textbf{99.2\%} \\
\bottomrule
\end{tabular}
\end{table}

\noindent\textbf{Provenance and language identification.} Recordings were curated from
publicly available film and classical vocal material and from regional folk or
traditional singing (forms such as \textit{thumri}, \textit{dadra}, \textit{chaiti} and
\textit{kajri}, and named regional traditions including Awadhi, Bhojpuri, Bundeli,
Garhwali, Kumaoni and Maithili). The language label is the language of the vocal content,
assigned during expert listening; code-mixed recordings with no dominant target language
were excluded.
% TODO(authors): per-recording source citation (archive / platform / URL / access date /
% licence) and the original sampling rate of each source file. Neither is recoverable
% from the release; both must come from the collection log.

\subsection{Released Clip Inventory}
Table~\ref{tab:s2} gives the clip count of every language$\times$partition cell of the
release, obtained by counting rows of the 24 \texttt{metadata.csv} files. These are the
counts a user obtains after download. Two cells are degenerate and should be treated as
absent for benchmarking purposes: Bengali So-VITS is empty (the folder exists but its
\texttt{metadata.csv} has no rows), and Punjabi So-VITS contains a single source
recording (21 clips).

\begin{table}[t]
\centering
\caption{Released clips per language and partition. Values are row counts of the
corresponding \texttt{metadata.csv}. All clips are exactly 10.0\,s.}
\label{tab:s2}
\footnotesize
\setlength{\tabcolsep}{4pt}
\begin{tabular}{lrrrrr}
\toprule
\textbf{Language} & \textbf{Bonafide} & \textbf{RVC} & \textbf{FFT} & \textbf{So-VITS} & \textbf{Total} \\
\midrule
Hindi   & 6{,}634 & 4{,}188 & 6{,}626 & 6{,}634 & 24{,}082 \\
Punjabi & 2{,}393 &   148 & 4{,}037 &    21 &  6{,}599 \\
Bengali & 2{,}806 & 2{,}806 & 2{,}806 &     0 &  8{,}418 \\
Telugu  & 4{,}489 &   103 & 4{,}245 & 4{,}481 & 13{,}318 \\
Odia    & 1{,}505 &    33 & 2{,}100 & 1{,}505 &  5{,}143 \\
Tamil   &   259 &    23 &   496 &   259 &  1{,}037 \\
\midrule
\textbf{Total} & \textbf{18{,}086} & \textbf{7{,}301} & \textbf{20{,}310} & \textbf{12{,}900} & \textbf{58{,}597} \\
\bottomrule
\end{tabular}
\end{table}

\subsection{Composition of the Synthetic Partitions}
\label{sec:composition}
The synthetic folders are not homogeneous: alongside conversion outputs they also carry
pipeline by-products that were swept into feature extraction and carry the label
\texttt{fake}. Table~\ref{tab:s3} resolves every synthetic cell into four content classes,
identified from the filename grammar of each entry (Section~\ref{sec:mapping}):

\begin{itemize}
  \item \emph{Converted} --- the conversion output proper.
  \item \emph{Vocal stem} --- files named \texttt{<title>\_vocals.wav}, i.e.\ isolated
        vocal tracks from the source-separation stage of the FFT branch.
  \item \emph{Reference} --- target-voice reference clips
        (\texttt{<lang>\_reference.wav}, \texttt{<lang>\_rvc\_ref.wav}).
  \item \emph{Unconverted source} --- in the Hindi RVC folder only, the 32 original
        recordings that sit beside their 32 \texttt{\_VOCAL\_CONVERTED} counterparts.
\end{itemize}

Users who require a strictly spoof-only synthetic partition should retain the
\emph{converted} column: \textbf{32{,}731} clips. The remaining 7{,}780 clips
(5{,}662 vocal stems, 24 references, 2{,}094 unconverted sources) are labelled
\texttt{fake} in the released metadata but are not conversion outputs, and are filtered by
the reference implementation in Section~\ref{sec:repro}.

\begin{table}[t]
\centering
\caption{Content classes inside the synthetic partitions (clips). Only the
\emph{Converted} column contains conversion outputs.}
\label{tab:s3}
\scriptsize
\setlength{\tabcolsep}{3pt}
\begin{tabular}{llrrrr}
\toprule
\textbf{Method} & \textbf{Lang.} & \textbf{Conv.} & \textbf{Stem} & \textbf{Ref.} & \textbf{Unconv.} \\
\midrule
\multirow{6}{*}{RVC}
 & Hindi   & 2{,}094 & --- & --- & 2{,}094 \\
 & Punjabi &   143 & --- &  5 & --- \\
 & Bengali & 2{,}806 & --- & --- & --- \\
 & Telugu  &    98 & --- &  5 & --- \\
 & Odia    &    28 & --- &  5 & --- \\
 & Tamil   &    18 & --- &  5 & --- \\
\midrule
\multirow{6}{*}{FFT}
 & Hindi   & 6{,}626 & --- & --- & --- \\
 & Punjabi & 2{,}004 & 2{,}032 &  1 & --- \\
 & Bengali & 2{,}806 & --- & --- & --- \\
 & Telugu  & 2{,}026 & 2{,}218 &  1 & --- \\
 & Odia    &   946 & 1{,}153 &  1 & --- \\
 & Tamil   &   236 &   259 &  1 & --- \\
\midrule
\multirow{6}{*}{So-VITS}
 & Hindi   & 6{,}634 & --- & --- & --- \\
 & Punjabi &    21 & --- & --- & --- \\
 & Bengali &     0 & --- & --- & --- \\
 & Telugu  & 4{,}481 & --- & --- & --- \\
 & Odia    & 1{,}505 & --- & --- & --- \\
 & Tamil   &   259 & --- & --- & --- \\
\midrule
\multicolumn{2}{l}{\textbf{Total}} & \textbf{32{,}731} & \textbf{5{,}662} & \textbf{24} & \textbf{2{,}094} \\
\bottomrule
\end{tabular}
\end{table}

\subsection{Singer and Voice Information}
\label{sec:singers}
Singer identity is recorded in the release only where the filename carries it, and no
gender field exists anywhere in the metadata. Concretely:

\noindent\textbf{Bonafide.} \texttt{original\_filename} holds the song title (Hindi,
Punjabi, Odia) or an index-prefixed title (Bengali, Telugu, Tamil). It does \emph{not}
name the performer. Bonafide singer identity is therefore not available from the release.
% TODO(authors): if a per-recording performer field exists in the collection log, add it
% as a `source_singer` column to Real_dataset/<lang>/metadata.csv and to the curation
% sheet. Gender, if required, should be annotated explicitly rather than inferred.

\noindent\textbf{RVC, Hindi.} This is the one branch whose filenames name voices. The 32
converted files are \texttt{<Voice>\_VOCAL\_CONVERTED.wav} over the following 32 named
voices: Akhil Sachdeva \& Tulsi Kumar; Alka Yagnik \& Arijit Singh; Allah Jilai Bai; Asha
Bhosle \& Mohammad Rafi; Begum Akhtar; Girija Devi; Kishore Kumar; Kishori Amonkar; Kumar
Sanu; Kumar Sanu \& Lata Mangeshkar; Lata Mangeshkar; Lata Mangeshkar \& Kishore Kumar;
M.S.\ Subbulakshmi; Maithili Thakur; Mohammad Rafi; Mukesh \& Lata Mangeshkar; Pandit
Bhimsen Joshi; Pandit Chhannulal Mishra; Pandit Jasraj; Pandit Kumar Gandharva; Roop Kumar
Rathod; Shobha Gurtu; Sonu Nigam \& Alka Yagnik; Udit Narayan; Udit Narayan \& Alka
Yagnik; Ustad Amir Khan; Ustad Rashid Khan; and five unnamed placeholders (Regional
Artist, Regional Folk Singer, Regional Singer, Traditional Artist, Traditional Singer).
One entry is truncated in the release as \texttt{M.S\_VOCAL\_CONVERTED.wav} (the source
name contains a period) and pairs with \texttt{M.S.~Subbulakshmi.wav}.
% TODO(authors): state whether these names identify the SOURCE performer or the TARGET
% voice model, and give the RVC model checkpoint used per row. The filename alone is
% ambiguous.

\noindent\textbf{RVC and FFT, other languages.} Filenames are sequence-numbered
(\texttt{rvc\_punjabi\_017.wav}) and carry no voice information.

\noindent\textbf{So-VITS.} Every So-VITS filename carries one of three fixed voice
presets rather than a target singer. Their distribution is given in
Table~\ref{tab:s4}; presets are assigned round-robin by sequence number, not chosen per
recording.

\begin{table}[t]
\centering
\caption{So-VITS voice presets: clips (source recordings) per language. Bengali is empty.}
\label{tab:s4}
\footnotesize
\begin{tabular}{lrrr}
\toprule
\textbf{Language} & \textbf{kawaii\_girl} & \textbf{soft\_anime} & \textbf{energetic\_idol} \\
\midrule
Hindi   & 2{,}894 (28) & 2{,}267 (28) & 1{,}473 (28) \\
Telugu  & 1{,}366 (17) & 1{,}591 (16) & 1{,}524 (17) \\
Odia    &   509 (5)  &   342 (5)  &   654 (5)  \\
Tamil   &    97 (4)  &    82 (3)  &    80 (3)  \\
Punjabi &    21 (1)  &     ---    &     ---    \\
\bottomrule
\end{tabular}
\end{table}

\subsection{Controlled Vocabularies}
The metadata uses closed vocabularies: \texttt{language} $\in$ \{\texttt{hindi},
\texttt{bengali}, \texttt{punjabi}, \texttt{tamil}, \texttt{telugu}, \texttt{odia}\};
\texttt{label} $\in$ \{\texttt{real}, \texttt{fake}\}; \texttt{method} $\in$
\{\texttt{real}, \texttt{rvc}, \texttt{fft}, \texttt{sovits}\}. Note the asymmetry
between the \texttt{method} value \texttt{sovits} and the directory name
\texttt{SoVits\_Anime}; both refer to the same partition.

\subsection{Train/Test Partitioning}
\label{sec:split}
Partitioning is performed at the \emph{recording} level, never the clip level, so that no
recording contributes clips to both sides and a detector cannot exploit
recording-specific cues. With seed 42, approximately 30\% of the recordings of each
language are held out, giving the assignment in Table~\ref{tab:s5}. The full song-level
assignment ships as \texttt{benchmark\_split.csv} (\texttt{SongID}, \texttt{Language},
\texttt{SongTitle}, \texttt{Clips}, \texttt{Split}), one row per bonafide source
recording.

\begin{table}[t]
\centering
\caption{Song-disjoint train/test split of the bonafide partition (seed 42). Clip counts
sum to the per-language totals of Table~\ref{tab:s1}.}
\label{tab:s5}
\footnotesize
\setlength{\tabcolsep}{4pt}
\begin{tabular}{lrrrr}
\toprule
\textbf{Language} & \textbf{Train songs} & \textbf{Test songs} & \textbf{Train clips} & \textbf{Test clips} \\
\midrule
Hindi   & 59 & 25 & 4{,}570 & 2{,}064 \\
Punjabi & 49 & 21 & 1{,}762 &   631 \\
Bengali & 45 & 19 & 2{,}030 &   776 \\
Telugu  & 35 & 15 & 2{,}131 & 2{,}358 \\
Odia    & 11 &  4 & 1{,}046 &   459 \\
Tamil   &  7 &  3 &   184 &    75 \\
\midrule
\textbf{Total} & \textbf{206} & \textbf{87} & \textbf{11{,}723} & \textbf{6{,}363} \\
\bottomrule
\end{tabular}
\end{table}

\noindent\textbf{Language-balanced test set.} From the held-out recordings, 60 bonafide
clips per language are drawn with seed 7, giving a balanced benchmark of
$60\times6=360$ clips. Sixty is bounded by Tamil, whose test partition offers 75 clips;
the margin is deliberately small so that no language is over-drained. The realised
manifest ships as \texttt{balanced\_test\_clips.csv} (\texttt{sample\_id},
\texttt{Language}, \texttt{original\_filename}, \texttt{clip\_index}). Tamil's 60
balanced clips derive from only 3 recordings, which limits the diversity --- not the
size --- of its test set; this is a property to report alongside Tamil results.

\noindent\textbf{Metrics.} Equal Error Rate is reported per language; any aggregate is
the equal-weight mean of per-language values, never a clip-weighted pool, so that Hindi
cannot dominate. 95\% confidence intervals are obtained by bootstrap resampling of the
balanced test clips with at least 1{,}000 resamples. The evaluation regime (zero-shot or
fine-tuned) is stated for every reported system.

% ============================================================================
\section{Feature-Extraction Specifications}
\label{sec:featspec}
% ============================================================================
All features are computed from the standardised 10\,s clips at a 16\,kHz sampling rate
(160{,}000 samples per clip), mono, peak-normalised. Table~\ref{tab:s6} gives the exact
released form of each feature: array shape, data type and the analysis settings implied
by that shape. Two distinct frame rates are used and this must be respected when features
are combined: the spectrogram-family features use a 25\,ms window with a 10\,ms hop
(\texttt{n\_fft}=400, \texttt{hop\_length}=160, giving $1+160000/160=1001$ frames), while
the pitch and flux features use librosa's default 2048-sample window with a 512-sample
hop (giving $1+\lfloor160000/512\rfloor=313$ frames).

\begin{table*}[t]
\centering
\caption{Released feature arrays. Shapes and dtypes were read from the files; settings
are those implied by the shapes and observed value ranges.}
\label{tab:s6}
\footnotesize
\setlength{\tabcolsep}{4pt}
\begin{tabular}{lllrp{7.4cm}}
\toprule
\textbf{Directory} & \textbf{Shape} & \textbf{dtype} & \textbf{Bytes/clip} & \textbf{Extraction settings} \\
\midrule
\texttt{logmel/} & $(80,1001)$ & float32 & 320{,}448 &
80 Mel bands, \texttt{n\_fft}=400 (25\,ms), \texttt{hop\_length}=160 (10\,ms),
\texttt{center}=True; power spectrogram converted to dB. Observed range
$[-59.95,\,20.05]$, consistent with \texttt{power\_to\_db(top\_db=80)}. \\
\texttt{mfcc/} & $(39,1001)$ & float32 & 156{,}284 &
13 MFCCs stacked with first- and second-order deltas ($13+13+13$), same framing as
\texttt{logmel}, computed from the 80-band log-Mel spectrogram. \\
\texttt{pitch\_frames/} & $(313,)$ & float64 & 2{,}632 &
Frame-wise $F_0$ by YIN. \texttt{frame\_length}=2048, \texttt{hop\_length}=512,
\texttt{center}=True. Values are bounded exactly by $[50,1000]$\,Hz, i.e.\
\texttt{fmin}=50, \texttt{fmax}=1000. Unvoiced frames are not flagged; the estimator
returns a value for every frame. \\
\texttt{spectral\_flux/} & $(312,)$ & float32 & 1{,}376 &
Frame-to-frame spectral difference on the same 2048/512 framing; length is one less than
the pitch frame count, as expected of a first difference. Non-negative. \\
\texttt{hnr\_frames/} & $(160000,)$ & float32 & 640{,}128 &
A per-\emph{sample} array (not per-frame) at the clip's full 16\,kHz resolution, whose
mean and standard deviation are released as \texttt{hnr\_mean} / \texttt{hnr\_std}.
Non-negative and heavy-tailed (observed max $3.1\times10^{5}$), i.e.\ a harmonic-to-noise
\emph{ratio} rather than a dB quantity. This array is 57\% of the download volume. \\
\texttt{whisper\_embeddings/} & $(512,)$ & float32 & 2{,}176 &
Whisper encoder hidden states mean-pooled over time. The 512-dimensional width
corresponds to the \texttt{base} Whisper encoder. \\
\bottomrule
\end{tabular}
\end{table*}
% TODO(authors): confirm from the extraction script (i) the exact HNR estimator and its
% units, (ii) the Whisper checkpoint and revision, and (iii) the librosa version. The
% released arrays fix the shapes and ranges but not the estimator internals.

Two derived quantities used in the paper are defined on these arrays. Jitter is the mean
absolute difference between consecutive pitch frames,
\begin{equation}
J = \frac{1}{N-1}\sum_{i=1}^{N-1}\left|F_0(i+1) - F_0(i)\right|,
\end{equation}
computed over the 313 frames of \texttt{pitch\_frames} and released as
\texttt{pitch\_jitter}; and the harmonic-to-noise ratio,
\begin{equation}
HNR = 10\log_{10}\!\left(\frac{E_{\text{harmonic}}}{E_{\text{noise}}}\right),
\end{equation}
with $E_{\text{harmonic}}$ and $E_{\text{noise}}$ the periodic and aperiodic energies.

Each folder's \texttt{summary\_features.csv} carries 13 clip-level scalars alongside
\texttt{sample\_id}, so that experiments needing only summary statistics require no
\texttt{.npy} download at all. The fields are \texttt{pitch\_mean},
\texttt{pitch\_variance}, \texttt{pitch\_jitter} (statistics of
\texttt{pitch\_frames}); \texttt{spectral\_flux\_mean}, \texttt{spectral\_flux\_std};
\texttt{hnr\_mean}, \texttt{hnr\_std}; \texttt{logmel\_mean}, \texttt{logmel\_std};
\texttt{mfcc\_mean}, \texttt{mfcc\_std} (grand mean and standard deviation over the whole
2-D array); and \texttt{whisper\_embedding\_norm},
\texttt{whisper\_embedding\_variance}. The \texttt{sample\_id} sets of
\texttt{metadata.csv} and \texttt{summary\_features.csv} are identical within every
folder, so the two join one-to-one.

% ============================================================================
\section{File Structure}
% ============================================================================

\subsection{Directory Layout}
\begin{scriptsize}
\begin{verbatim}
MSDF-India/
|- Real_dataset/
|  |- hindi/          bengali/   punjabi/
|  |- tamil/          telugu/    odia/
|     |- logmel/               *.npy
|     |- mfcc/                 *.npy
|     |- pitch_frames/         *.npy
|     |- spectral_flux/        *.npy
|     |- hnr_frames/           *.npy
|     |- whisper_embeddings/   *.npy
|     |- metadata.csv
|     |- summary_features.csv
|- Deepfake_dataset/
   |- RVC/
   |  |- hindi_rvc/  bengali_rvc/  punjabi_rvc/
   |  |- tamil_rvc/  telugu_rvc/   odia_rvc/
   |     |- (same six feature dirs + 2 CSVs)
   |- FFT/
   |  |- <lang>_fft/  ...
   |- SoVits_Anime/
      |- <lang>_sovits_anime/  ...
\end{verbatim}
\end{scriptsize}

\noindent The layout is uniform: 24 leaf folders (6 bonafide + 18 synthetic), each with
the same six feature directories and the same two CSV files, and each self-describing.
A user interested in a single language$\times$method cell needs to download exactly one
leaf folder.

\subsection{Feature-File Naming Convention}
\label{sec:naming}
Feature files follow
\begin{center}
\texttt{<prefix>\_<stem>\_clip<k>.npy}
\end{center}
where \texttt{<stem>} is the value of \texttt{original\_filename} with its final
extension removed, \texttt{<k>} is \texttt{clip\_index} (0-based, no zero padding), and
\texttt{<prefix>} identifies the directory the source audio was read from during
extraction. The same file name is used in all six feature directories of a folder, so
one clip is addressed by one name across every feature type.

Observed prefixes are \texttt{<language>\_real} for the bonafide partition,
\texttt{<language>\_rvc} and \texttt{<language>\_fft} for those branches,
\texttt{<language>\_sovits\_anime} or \texttt{output\_sovits\_anime} for So-VITS
(the prefix differs by language), and \texttt{references\_rvc} /
\texttt{references\_fft} for the target-voice reference clips.

\noindent\textbf{Practical consequence.} The prefix is \emph{not} a column of
\texttt{metadata.csv}, so a feature path cannot be constructed from a metadata row alone.
The robust join is by suffix:
\begin{scriptsize}
\begin{verbatim}
suffix = f"_{stem}_clip{row.clip_index}.npy"
path   = next(p for p in listdir(featdir)
              if p.endswith(suffix))
\end{verbatim}
\end{scriptsize}
\noindent which is what the reference loader in Section~\ref{sec:repro} does.

\subsection{Metadata-Field Definitions}
Every leaf folder's \texttt{metadata.csv} has exactly the seven fields of
Table~\ref{tab:s7}, one row per released clip, and no missing values.

\begin{table*}[t]
\centering
\caption{\texttt{metadata.csv} field definitions. The same seven fields appear in all 24
leaf folders, with no missing values.}
\label{tab:s7}
\footnotesize
\begin{tabular}{lp{14.2cm}}
\toprule
\textbf{Field} & \textbf{Definition} \\
\midrule
\texttt{sample\_id} & 32-character lowercase hexadecimal clip identifier. Unique across
the entire release (58{,}597 rows, 58{,}597 distinct values). Opaque: it is not a hash of
any released field, and must be treated as a primary key rather than recomputed. It is
the join key to \texttt{summary\_features.csv} and to
\texttt{balanced\_test\_clips.csv}. \\
\texttt{language} & Language code. In 12 of the 24 files this column instead carries the
source audio's parent directory name---\texttt{vocals} or \texttt{references} in the FFT and
RVC folders of Odia, Punjabi, Tamil and Telugu, and \texttt{output} for \emph{every} So-VITS
row of those four languages (11{,}952 clips in total). The containing folder is
authoritative; grouping on this column as released silently drops 6{,}266 So-VITS clips into
a spurious ``output'' language. \\
\texttt{label} & \texttt{real} for bonafide clips, \texttt{fake} for every clip in a
synthetic folder --- including the non-conversion content of
Section~\ref{sec:composition}. \\
\texttt{method} & Generation method: \texttt{real}, \texttt{rvc}, \texttt{fft} or
\texttt{sovits}. Constant within a folder. \\
\texttt{original\_filename} & File name of the audio the clip was cut from, including its
extension. This is the only pointer back to a source recording, and the key used for the
mapping of Section~\ref{sec:mapping}. \\
\texttt{clip\_index} & 0-based index of the clip inside its source recording. Contiguous:
a recording of $n$ clips uses indices $0\ldots n-1$. \\
\texttt{duration\_seconds} & Clip duration in seconds. Always 10 across all 58{,}597
rows. \\
\bottomrule
\end{tabular}
\end{table*}

\subsection{An Example Feature File}
Fig.~\ref{fig:example} works one clip end to end: its metadata row, the six arrays it
resolves to, and its summary-feature row.

\begin{figure*}[t]
\centering
\begin{scriptsize}
\begin{verbatim}
# Real_dataset/tamil/metadata.csv, row 1
sample_id,language,label,method,original_filename,clip_index,duration_seconds
f3bbf40fcea681d7d15dbad01f0d6665,tamil,real,real,004_Hello Brother.mp3.mp3,0,10

# The six arrays this row resolves to (identical basename in every feature directory):
Real_dataset/tamil/logmel/tamil_real_004_Hello Brother.mp3_clip0.npy      (80, 1001) float32
Real_dataset/tamil/mfcc/tamil_real_004_Hello Brother.mp3_clip0.npy        (39, 1001) float32
Real_dataset/tamil/pitch_frames/...  _clip0.npy                           (313,)     float64
Real_dataset/tamil/spectral_flux/... _clip0.npy                           (312,)     float32
Real_dataset/tamil/hnr_frames/...    _clip0.npy                           (160000,)  float32
Real_dataset/tamil/whisper_embeddings/... _clip0.npy                      (512,)     float32

# Real_dataset/tamil/summary_features.csv, matching row
sample_id            = f3bbf40fcea681d7d15dbad01f0d6665
pitch_mean           = 136.994      pitch_variance      = 29617.367   pitch_jitter  = 17.463
spectral_flux_mean   = 9.161        spectral_flux_std   = 11.775
hnr_mean             = 39.481       hnr_std             = 1312.259
logmel_mean          = -45.702      logmel_std          = 18.236
mfcc_mean            = -5.177       mfcc_std            = 78.724
whisper_embedding_norm = 21.927     whisper_embedding_variance = 0.939

# Note: original_filename ends in ".mp3.mp3"; the stem used in the file name is
# "004_Hello Brother.mp3" (only the FINAL extension is stripped). File names contain
# spaces and periods, so paths must be quoted.
\end{verbatim}
\end{scriptsize}
\caption{One released clip, from metadata row to feature arrays to summary statistics.
Values are verbatim from the release.}
\label{fig:example}
\end{figure*}

\subsection{Mapping Feature Files to Source Recordings}
\label{sec:mapping}
Because the release is feature-level, the link between a synthetic clip and the bonafide
recording it was derived from lives entirely in \texttt{original\_filename}. The filename
grammars differ by language and method, and so does the strength of the link:

\begin{itemize}
  \item \textbf{Bonafide:} \texttt{<Title>.mp3} (Hindi, Punjabi, Odia) or
        \texttt{<NNN>\_<Title>.mp3.mp3} (Bengali, Telugu, Tamil), where \texttt{<NNN>} is
        a per-language source index.
  \item \textbf{So-VITS, all languages:}
        \texttt{anime\_<preset>\_<seq>\_<Title>.wav} (Hindi, Odia) or
        \texttt{anime\_<preset>\_<seq>\_<NNN>\_<Title>.mp3.wav} (Punjabi, Tamil,
        Telugu). Both carry the title, so the mapping is exact.
  \item \textbf{FFT, Hindi:} \texttt{deepfake\_<seq>\_<srcid>\_<Title>.mp3} --- the only
        branch with an explicit numeric source pointer in addition to the title.
  \item \textbf{FFT and RVC, Bengali:} \texttt{fft\_<seq>\_<Title>.mp3\_22k.wav} and
        \texttt{rvc\_fixed\_<seq>\_<Title>.mp3\_22k.wav}; title-carrying, exact.
  \item \textbf{FFT, other languages:} \texttt{fft\_<lang>\_<seq>.wav} for the conversion
        outputs (no title) alongside \texttt{<Title>\_vocals.wav} for the vocal stems
        (title-carrying).
  \item \textbf{RVC, other languages:} \texttt{rvc\_<lang>\_<seq>.wav}; sequence number
        only.
  \item \textbf{RVC, Hindi:} \texttt{<Voice>.wav} and
        \texttt{<Voice>\_VOCAL\_CONVERTED.wav}; a voice-name namespace disjoint from the
        bonafide title namespace.
\end{itemize}

Table~\ref{tab:s8} quantifies the resulting coverage: the fraction of synthetic source
files in each cell that can be matched to a bonafide source recording by
normalised-title matching. Coverage is complete for every So-VITS cell, for Hindi FFT and
for both Bengali branches; it is 50\% for the FFT cells of Odia, Punjabi, Tamil and
Telugu (exactly the vocal-stem half); and it is zero for RVC outside Bengali, where the
filenames carry only a sequence number.

\begin{table}[t]
\centering
\caption{Mapping coverage: synthetic source files resolvable to a bonafide source
recording by title, per cell.}
\label{tab:s8}
\footnotesize
\setlength{\tabcolsep}{4pt}
\begin{tabular}{lrrr}
\toprule
\textbf{Language} & \textbf{RVC} & \textbf{FFT} & \textbf{So-VITS} \\
\midrule
Hindi   & 0/64 (0\%)   & 84/84 (100\%)  & 84/84 (100\%) \\
Bengali & 64/64 (100\%)& 64/64 (100\%)  & --- (empty) \\
Punjabi & 0/71 (0\%)   & 69/138 (50\%)  & 1/1 (100\%) \\
Telugu  & 0/51 (0\%)   & 45/90 (50\%)   & 50/50 (100\%) \\
Odia    & 0/16 (0\%)   & 14/28 (50\%)   & 15/15 (100\%) \\
Tamil   & 0/11 (0\%)   & 10/20 (50\%)   & 10/10 (100\%) \\
\bottomrule
\end{tabular}
\end{table}

% TODO(authors): release a `source_map.csv` with columns
% (method, language, original_filename, source_song_id, source_original_filename)
% built from the generation logs. That is the one artefact that would make the
% sequence-numbered RVC and FFT branches joinable; it cannot be reconstructed from the
% released names.

Song-disjoint splitting of the \emph{synthetic} partitions is therefore currently exact
only for the cells with full coverage. For the sequence-numbered cells, the reference
implementation splits at the level of \texttt{original\_filename}, which still guarantees
that no synthetic recording contributes clips to both partitions, but cannot guarantee
that a synthetic clip and its bonafide counterpart land on the same side. This is the
main reason the headline benchmark of the paper is defined on the bonafide-balanced test
set.

% ============================================================================
\section{Reproducing the Benchmark}
\label{sec:repro}
% ============================================================================
The benchmark is reproducible from the released CSVs alone; the \texttt{.npy} arrays are
needed only for systems that consume frame-level features.

\begin{enumerate}
  \item \textbf{Obtain the release.} Download the Drive folder (or only the leaf folders
        of interest) preserving the directory layout.
  \item \textbf{Build the clip index.} Concatenate the 24 \texttt{metadata.csv} files,
        adding the leaf folder path as an extra column. This yields 58{,}597 rows.
  \item \textbf{Filter the synthetic partition} to conversion outputs only, dropping
        names ending in \texttt{\_vocals.wav}, names matching
        \texttt{*reference*}/\texttt{*\_ref*}, and, in Hindi RVC, names without the
        \texttt{\_VOCAL\_CONVERTED} marker. This leaves 32{,}731 synthetic clips
        (Table~\ref{tab:s3}).
  \item \textbf{Apply the fixed split.} Join \texttt{benchmark\_split.csv} on the
        bonafide \texttt{original\_filename} to obtain train/test at the recording level
        (Table~\ref{tab:s5}). Do not re-derive the split with a fresh seed; the released
        assignment is the protocol.
  \item \textbf{Load the balanced test set.} Read
        \texttt{balanced\_test\_clips.csv} and select those 360 \texttt{sample\_id}
        values, adding the matching synthetic clips per language and per method.
  \item \textbf{Score and report.} Compute EER per language, aggregate as the
        equal-weight mean over languages, and attach bootstrap 95\% confidence intervals
        ($\geq$1{,}000 resamples) to every per-language value. State whether the system
        was evaluated zero-shot or fine-tuned.
\end{enumerate}

Fig.~\ref{fig:code} gives a minimal reference loader implementing steps 2, 3 and the
suffix join of Section~\ref{sec:naming}.

\begin{figure*}[t]
\centering
\begin{scriptsize}
\begin{verbatim}
import os, re, glob, numpy as np, pandas as pd

ROOT = "MSDF-India"

def load_index(root=ROOT):
    """Concatenate all 24 metadata.csv files into one clip index."""
    rows = []
    for csv in glob.glob(os.path.join(root, "*_dataset", "**", "metadata.csv"), recursive=True):
        d = pd.read_csv(csv)
        if len(d):
            d["folder"] = os.path.dirname(csv)
            rows.append(d)
    return pd.concat(rows, ignore_index=True)          # 58,597 rows

def conversion_outputs_only(df):
    """Drop vocal stems, target-voice references and unconverted Hindi RVC sources."""
    fn   = df.original_filename
    stem = fn.str.lower().str.endswith("_vocals.wav")
    ref  = fn.str.lower().str.contains(r"reference|_ref\b|_ref\.", regex=True)
    unc  = (df.method.eq("rvc") & df.language.eq("hindi")
            & ~fn.str.contains("_VOCAL_CONVERTED"))
    return df[~(stem | ref | unc)]                     # 18,086 real + 32,731 fake

def feature_path(row, feature="logmel"):
    """Resolve one metadata row to its .npy file (the folder prefix is not in the CSV)."""
    stem   = row.original_filename.rsplit(".", 1)[0]
    suffix = f"_{stem}_clip{row.clip_index}.npy"
    d      = os.path.join(row.folder, feature)
    for p in os.listdir(d):
        if p.endswith(suffix):
            return os.path.join(d, p)
    raise FileNotFoundError(suffix)

idx   = conversion_outputs_only(load_index())
split = pd.read_csv("benchmark_split.csv")             # SongID,Language,SongTitle,Clips,Split
bal   = pd.read_csv("balanced_test_clips.csv")         # 360 balanced bonafide sample_ids
test  = idx[idx.sample_id.isin(set(bal.sample_id))]

x = np.load(feature_path(idx.iloc[0]))                 # (80, 1001) float32
\end{verbatim}
\end{scriptsize}
\caption{Reference loader: builds the clip index, filters the synthetic partition to
conversion outputs, and resolves a metadata row to its feature array.}
\label{fig:code}
\end{figure*}

\noindent\textbf{Environment.} The released arrays are plain \texttt{numpy} \texttt{.npy}
files and load with any recent \texttt{numpy}; \texttt{pandas} is the only other
requirement for the CSVs. Re-extracting features from audio additionally requires
\texttt{librosa} and, for the Whisper embeddings, the \texttt{transformers} Whisper
encoder.
% TODO(authors): pin the exact versions used for the released extraction
% (python, numpy, librosa, torch, transformers) so re-extraction is bit-comparable.

% ============================================================================
\section{Code Repository and README}
% ============================================================================
The code accompanying the release is organised as follows:

\begin{scriptsize}
\begin{verbatim}
msdf-india/
|- README.md
|- requirements.txt
|- data/
|  |- benchmark_split.csv
|  |- balanced_test_clips.csv
|  |- selection_checklist.xlsx
|- src/
|  |- extract_features.py  # audio -> npy + CSV
|  |- build_index.py       # merge 24 metadata.csv
|  |- filter_synthetic.py  # conversion outputs
|  |- make_splits.py       # seeds 42 and 7
|  |- evaluate.py          # EER + bootstrap CIs
|- notebooks/
   |- selection_validation.ipynb  # PCA/kmeans
\end{verbatim}
\end{scriptsize}

\noindent The repository URL is
% TODO(authors): insert the public repository URL here before submission, and mirror
% benchmark_split.csv / balanced_test_clips.csv into the Drive release so the protocol
% travels with the data.
\texttt{[to be inserted]}.

\noindent The \texttt{README} states, in order: the release URL and licence terms; the
directory layout of Section V-A; the field definitions of Table~\ref{tab:s7}; the feature
specification of Table~\ref{tab:s6}; the six-step reproduction procedure of
Section~\ref{sec:repro}; the filtering rule of Section~\ref{sec:composition} together
with the resulting clip counts; the known issues of Section~\ref{sec:issues}; and the
citation and contact details for requesting the original audio under the controlled-use
agreement.

% ============================================================================
\section{Release Notes and Known Issues}
\label{sec:issues}
% ============================================================================
The following properties of the current release are documented so that users are not
surprised by them, and so that any number computed from the archive can be reconciled
with the paper.

\begin{enumerate}
  \item \textbf{Non-conversion content in the synthetic folders.} 7{,}780 clips labelled
        \texttt{fake} are vocal stems, target-voice references, or (in Hindi RVC) the
        unconverted source recordings themselves. They must be filtered before training
        or evaluation; Section~\ref{sec:repro} does so. Treating the synthetic folders as
        wholly spoof inflates the spoof partition by 19\%.
  \item \textbf{Language field.} The \texttt{language} column is unreliable as released
        (Table~\ref{tab:s7}): in half the files it carries the extraction directory name
        rather than the language. Every language-wise figure in this document is derived
        from the containing folder, which is authoritative, and the loader released with the
        code repository repairs the column on load.
        % TODO(authors): regenerate the 12 affected metadata.csv files with the correct
        % language value before the next upload.
  \item \textbf{Degenerate cells.} Bengali So-VITS is empty and Punjabi So-VITS contains
        one source recording. The So-VITS branch is therefore evaluable on four languages,
        not six.
  \item \textbf{Directory and label naming.} The So-VITS partition is stored under
        \texttt{SoVits\_Anime/} and its files carry one of three fixed voice presets
        (Table~\ref{tab:s4}) rather than a per-recording target singer, while
        \texttt{metadata.csv} records \texttt{method=sovits}.
        % TODO(authors): this partition must be reconciled with the So-VITS description in
        % the main paper before submission -- either by regenerating it with the neural
        % pipeline and re-extracting features, or by describing the released variant. Do
        % not leave the released artefacts and the paper text describing different systems.
  \item \textbf{Counts.} The clip counts of Table~\ref{tab:s2} are the authoritative
        description of the release. They are computed at the clip level and are not
        directly comparable with recording-level file counts and durations reported for
        the generation stage.
        % TODO(authors): reconcile Table III of the main paper (recording-level file
        % counts of the generation store, e.g. Hindi RVC = 241) with the released feature
        % counts here (Hindi RVC = 64 source files / 4,188 clips). Either report both and
        % label them, or report only the released population.
  \item \textbf{Storage asymmetry.} \texttt{hnr\_frames} stores one value per audio sample
        and accounts for 57\% of the archive. A future release could store it at the same
        frame rate as the other descriptors at negligible information cost.
  \item \textbf{Identifiers.} \texttt{sample\_id} is opaque and cannot be recomputed from
        the released fields; it must be carried through any derived artefact so that
        results remain joinable to the release.
\end{enumerate}

\end{document}
